\chapter{Debugging Real Failures}\label{ch:gdb-failures}

The previous chapters were the instrument panel. This one is the flying: the
handful of failures you will actually meet, and the shortest route through each.

\section{Segmentation Faults}\label{sec:gdb-segv}

\begin{definition}[Segmentation fault]
	The kernel's response to a memory access the process is not allowed to make:
	dereferencing a null or wild pointer, running off the end of the stack, or
	writing to read-only memory. The signal is \texttt{SIGSEGV}
	(\autoref{sec:sh-signals}).
\end{definition}

The procedure is mechanical:

\begin{gdbcode}
(gdb) run
Program received signal SIGSEGV, Segmentation fault.
0x00005555555551a9 in parse_line (s=0x0) at parser.c:42
42          while (*s != '\n') {
(gdb) bt                        # 1. how did I get here?
#0  parse_line (s=0x0) at parser.c:42
#1  0x0000555555555241 in read_config (path=0x7fff...) at config.c:18
#2  0x00005555555552c8 in main (argc=2, argv=0x7fff...) at main.c:11
(gdb) info args                 # 2. what was I given?
s = 0x0
(gdb) frame 1                   # 3. who gave it to me?
#1  read_config (path=0x7fffffffe5c1 "app.conf") at config.c:18
18          parse_line(next_line(f));
(gdb) print f
$1 = (FILE *) 0x0               # 4. and why? fopen failed, unchecked
\end{gdbcode}

\begin{moral}
	\cmd{bt}, \cmd{info args}, \cmd{up} until something looks wrong. Four commands
	solve the large majority of segfaults, and none of them requires understanding
	the program first.
\end{moral}

\begin{keys}[0.30]
	\texttt{0x0}                 & A null pointer: something returned \texttt{NULL} unchecked \\
	A small value like \texttt{0x10} & A null pointer plus a struct offset --- \texttt{p->field} where \texttt{p} is null \\
	\texttt{0xffffffffffffffff}  & An uninitialised or already-freed pointer          \\
	An address just past a valid one & A buffer overrun                               \\
	The address is fine, the write fails & Writing to a string literal or other read-only memory \\
\end{keys}

\begin{note}
	When the pointer expression is complicated, GDB will tell you the faulting
	address directly:
\begin{gdbcode}
(gdb) print $_siginfo._sifields._sigfault.si_addr
$1 = (void *) 0x0
\end{gdbcode}
\end{note}

\section{Stack Overflow and Runaway Recursion}\label{sec:gdb-stackoverflow}

A crash whose backtrace is thousands of identical frames, or entirely
\texttt{??}, is usually the stack rather than the heap.

\begin{gdbcode}
(gdb) bt
#0  descend (n=0x5555555592a0) at tree.c:14
#1  0x0000555555555199 in descend (n=0x5555555592a0) at tree.c:17
#2  0x0000555555555199 in descend (n=0x5555555592a0) at tree.c:17
...
#65531 0x0000555555555199 in descend (n=0x5555555592a0) at tree.c:17
(gdb) bt -5                    # the OUTERMOST frames: where it started
(gdb) frame 65000
(gdb) print n == $_saved        # same node every time -- a cycle in the tree
\end{gdbcode}

\begin{keys}[0.30]
	\cmd{bt 10}, \cmd{bt -10}   & The two ends of the stack; the middle is repetition \\
	\cmd{set backtrace limit 20} & Stop GDB printing 60,000 frames                    \\
	\cmd{info frame}            & Where the stack pointer has reached                 \\
	\cmd{ulimit -s}             & The stack limit, usually 8\,MB                      \\
\end{keys}

\begin{note}
	Infinite recursion and a merely deep recursion look identical in a backtrace.
	The distinguishing question is whether the argument changes between frames: if
	frame \texttt{\#1} and frame \texttt{\#500} have the same pointer, the data
	structure has a cycle. A very large stack array --- \cmd{char buf[8*1024*1024]}
	--- produces the same crash with a backtrace only one frame deep.
\end{note}

\section{Uninitialised Memory and Heap Corruption}\label{sec:gdb-memory}

These are the failures GDB is \emph{worst} at, and it is important to know why.

\begin{intuition}
	A use-after-free or a buffer overrun does its damage at one moment and causes a
	crash at another, often far away and in unrelated code. GDB shows you the
	crash. It has no way to show you the write that caused it, because by then
	nothing is left to see.
\end{intuition}

What GDB can still do:

\begin{keys}[0.30]
	\cmd{watch -l *ptr}      & Catch the write, if you can guess the address in advance \\
	\cmd{x/32xb p-16}        & Look at the bytes around a suspect block                 \\
	\cmd{print *p}           & Implausible field values suggest the block was reused    \\
	\cmd{p \$\_memeq(a, b, n)} & Compare two regions                                    \\
\end{keys}

\begin{gdbcode}
(gdb) print *node
$1 = {name = 0x6161616161616161 <error: Cannot access memory>, port = 1633771873}
\end{gdbcode}

\begin{note}
	\texttt{0x6161616161616161} is \texttt{"aaaaaaaa"} read as a pointer --- a
	string has been written over a struct. Recognising ASCII in a pointer is a
	useful reflex: \texttt{0x41414141} is \texttt{"AAAA"},
	\texttt{0xdeadbeef} and \texttt{0xcdcdcdcd} are debug fill patterns, and
	\texttt{0x0} is the honest one.
\end{note}

\begin{moral}
	Do not hunt memory corruption in GDB. Rebuild with AddressSanitizer, which
	catches the bad access at the moment it happens and prints both the access and
	the allocation it belongs to. Then, if you still need to, come back to GDB
	knowing exactly where to put the breakpoint.
\end{moral}

\begin{shell}
$ gcc -g -O0 -fsanitize=address,undefined -o app main.c
$ ./app
==12345==ERROR: AddressSanitizer: heap-use-after-free on address 0x60300000eff4
READ of size 4 at 0x60300000eff4 thread T0
    #0 0x55d in use_node table.c:88
freed by thread T0 here:
    #1 0x4f2 in release table.c:61
previously allocated by thread T0 here:
    #2 0x48a in make_node table.c:22
\end{shell}

\begin{note}
	Three stack traces --- the bad access, the \cmd{free}, and the
	\cmd{malloc} --- which is the entire life story of the bug. This is available
	on this machine (\autoref{sec:gdb-companions}) and costs one compiler flag.
\end{note}

\section{Multi-Threaded Programs}\label{sec:gdb-threads}

\begin{keys}[0.32]
	\cmd{info threads}              & List them; \texttt{*} marks the current one    \\
	\cmd{thread 3}                  & Switch to thread 3                             \\
	\cmd{thread apply all bt}       & A backtrace of every thread                    \\
	\cmd{thread apply all bt full}  & \dots{} with locals. The deadlock command      \\
	\cmd{break f thread 3}          & A breakpoint only thread 3 honours             \\
	\cmd{set scheduler-locking on}  & Step one thread; freeze the others             \\
	\cmd{set non-stop on}           & Stop one thread while the rest keep running    \\
\end{keys}

\begin{eg}[Diagnosing a deadlock]
	Attach to the hung process (\autoref{sec:gdb-start}) and ask every thread what
	it is doing:
\begin{gdbcode}
(gdb) thread apply all bt

Thread 2 (Thread 0x7ffff7a3d700):
#0  futex_wait (...) at lowlevellock.c:52
#1  __lll_lock_wait (...) 
#2  work (arg=0x0) at worker.c:44        # waiting for mutex B
    
Thread 1 (Thread 0x7ffff7fc0740):
#0  futex_wait (...) at lowlevellock.c:52
#2  main () at main.c:31                 # waiting for mutex A
\end{gdbcode}
	Two threads, each blocked in \cmd{futex\_wait}, each holding what the other
	wants: a lock-ordering deadlock. One command diagnosed it.
\end{eg}

\begin{note}
	\cmd{set scheduler-locking on} makes stepping comprehensible --- otherwise
	other threads run while you step and the state changes under you. It also makes
	deadlock easy to cause yourself, so turn it off again before \cmd{continue}.
\end{note}

\begin{moral}
	\cmd{thread apply all bt} is to concurrency what \cmd{bt} is to a crash. For a
	hang, it is very often the whole diagnosis.
\end{moral}

\section{Handling Signals}\label{sec:gdb-signals}

GDB intercepts every signal before the program sees it, and you can change what
it does with each.

\begin{keys}[0.32]
	\cmd{info signals}                & The whole table                              \\
	\cmd{handle SIGPIPE nostop noprint} & Ignore it entirely                          \\
	\cmd{handle SIGUSR1 stop print}   & Stop on it                                    \\
	\cmd{handle SIGSEGV nostop pass}  & Let the program's own handler have it         \\
	\cmd{signal SIGINT}               & Deliver a signal by hand                      \\
	\cmd{continue}                    & Resume, delivering the pending signal         \\
\end{keys}

\begin{note}
	The common need is \cmd{handle SIGPIPE nostop noprint pass}: a networking
	program may take \texttt{SIGPIPE} routinely, and GDB stopping every time makes
	the session unusable. The opposite case is a program with its own
	\texttt{SIGSEGV} handler --- GDB stops first, which is what you want, and
	\cmd{pass} lets the handler run if you want to debug the handler instead.
\end{note}

\section{Companion Tools}\label{sec:gdb-companions}

GDB is an inspector, not a detector. These find bugs it can only witness.

\begin{keys}[0.30]
	\vocab{AddressSanitizer} & \cmd{-fsanitize=address}. Use-after-free, buffer overruns, leaks, with the allocating stack trace. A few times slower \\
	\vocab{UBSan}            & \cmd{-fsanitize=undefined}. Signed overflow, bad shifts, misaligned access, null dereference \\
	\vocab{ThreadSanitizer}  & \cmd{-fsanitize=thread}. Data races, which are otherwise nearly impossible to find \\
	\vocab{Valgrind}         & \cmd{valgrind ./app}. The same class of errors with no recompilation, considerably slower \\
	\vocab{rr}               & \cmd{rr record} then \cmd{rr replay}. Records an execution so you can step \emph{backwards} in GDB \\
	\vocab{strace}, \vocab{ltrace} & Every system or library call the program makes \\
\end{keys}

\begin{note}
	On this machine the sanitisers are available through GCC 13, and \cmd{strace}
	is present; \cmd{valgrind} and \cmd{rr} are not installed
	(\cmd{apt install valgrind}). The sanitisers are the better default anyway ---
	they are faster and more precise --- with Valgrind's advantage being that it
	needs no rebuild, which matters when you only have the binary.
\end{note}

\begin{moral}
	Match the tool to the failure. A crash with a sensible backtrace: GDB. Memory
	corruption, a leak, or a crash whose backtrace is nonsense: a sanitiser first,
	then GDB. A hang: \cmd{gdb -p} and \cmd{thread apply all bt}. A wrong answer
	with no crash at all: a breakpoint and \cmd{display}.
\end{moral}
